\documentclass[11pt]{article}
\usepackage[margin=1in]{geometry}
\usepackage{amsmath,amssymb,amsthm}
\usepackage[hidelinks]{hyperref}
\emergencystretch=2em
\newtheorem{theorem}{Theorem}
\newtheorem{lemma}{Lemma}
\newcommand{\Pp}{\mathbb P}
\newcommand{\Gstar}{\mathcal G_*}
\title{Counterexample to Conjecture 9305:\\Local Weak Limits Need Not Be Supported on Transitive Graphs}
\author{gaochengzhecpu}
\date{}
\begin{document}
\maketitle

\paragraph{The claim being disproved.}
The source asserts that local limits admit a decomposition into extreme
points whose supports consist of transitive graphs.  Here ``transitive''
means vertex-transitive: the automorphism group acts transitively on the
vertices of the underlying graph.  We disprove the support assertion.
In fact the counterexample cannot be a mixture of \emph{any} laws
concentrated on vertex-transitive graphs, regardless of extremality.
The source places no connectedness or infinitude requirement on the
approximating finite graphs or their limiting root components.

\paragraph{The rooted graph space.}
Let $\Gstar$ consist of rooted isomorphism classes of countable, connected,
locally finite simple graphs.  The radius-$r$ ball is the induced graph on
vertices joined to the root by a walk of length at most $r$; isomorphisms
of balls must preserve the root.  Equip $\Gstar$ with its local sigma
algebra generated by these ball-isomorphism cylinders.  The uniformly
rooted law of a finite graph means: choose one of its actual vertices
uniformly, take its actual connected component, and retain the root.
Thus disconnected approximating graphs cause no ambiguity. This is the
standard uniform-root convention of Aldous and Lyons~\cite{AL},\S2, p.~10.

Let $P_3$ be the path $0$--$1$--$2$.  Write $[P_3,r]$ for its rooted
isomorphism class, and define the probability measure
\[
 \mu=\frac13\sum_{r=0}^2\delta_{[P_3,r]}.
\]
Its two distinct rooted classes have masses $2/3$ (an endpoint root)
and $1/3$ (the middle root).

\begin{lemma}[A genuine local weak limit]
Let $G_m$ be the disjoint union of $m+1$ copies of $P_3$.  Then
$|V(G_m)|=3(m+1)\to\infty$, and the uniformly rooted laws of $G_m$
are exactly $\mu$ for every $m\geq0$.  In particular they converge
locally weakly to $\mu$.
\end{lemma}
\begin{proof}
Use the actual vertex set
$\{0,\ldots,m\}\times\{0,1,2\}$, with adjacency precisely when the
first coordinates agree and the second coordinates are adjacent in
$P_3$.  The component of $(i,r)$ is the $i$th copy, and its rooted
isomorphism class is $[P_3,r]$.  Each second coordinate occurs $m+1$
times among the $3(m+1)$ vertices.  Consequently the pushforward of
the uniform vertex law is exactly the displayed measure.  Every
rooted ball cylinder therefore has the same probability for every $m$;
this is the standard cylinder criterion for local weak convergence.
\end{proof}

\newpage
\paragraph{A separating cylinder event.}
For a finite rooted graph $(H,o)$, let $C_r(H,o)$ be the event that the
rooted radius-$r$ ball is isomorphic to that of $(H,o)$.  Set
\[
 A=\bigcup_{r=0}^2 C_2(P_3,r).
\]
This is a finite union of genuine local cylinders, hence measurable,
and $\mu(A)=1$.

\begin{lemma}[No vertex-transitive graph belongs to $A$]
For every countable connected locally finite vertex-transitive graph
$(H,o)$, one has $[H,o]\notin A$.
\end{lemma}
\begin{proof}
All vertices of a vertex-transitive locally finite graph have the same
degree, because any automorphism induces a bijection between the
corresponding neighbor sets.  Suppose the rooted radius-two ball at $o$
is a rooted copy of $P_3$.  Every neighbor of $o$, and every neighbor
of such a neighbor, belongs to this radius-two ball.  Therefore the
degree of the root and the degree of each of its neighbors can both
be read from the induced ball: neither neighbor set is truncated.

If the root corresponds to an endpoint of $P_3$, its degree is $1$
and its neighbor's degree is $2$.  If it corresponds to the middle
vertex, its degree is $2$ and each neighbor's degree is $1$.  Either
case contradicts equality of degrees in a vertex-transitive graph.
This argument does not assert that a radius-two $P_3$ ball identifies
the whole component; for example an endpoint of $P_4$ has such a ball.
\end{proof}

\newpage
\begin{theorem}[Failure of the claimed decomposition]
There are no mixing probability measure $\eta$ and measurable family
of probability measures $(\nu_i)$ on $\Gstar$ such that
\[
 \mu(B)=\int \nu_i(B)\,d\eta(i)
 \quad\hbox{for every measurable } B,
\]
and $\nu_i$ is concentrated on vertex-transitive graphs for
$\eta$-almost every $i$.
\end{theorem}
\begin{proof}
By the separating-cylinder lemma, concentration on vertex-transitive
graphs implies $\nu_i(A)=0$.  Substituting the measurable event $A$
in the barycentric identity gives
\[
 1=\mu(A)=\int \nu_i(A)\,d\eta(i)=0,
\]
a contradiction.  The argument even allows arbitrary measures for
the components and the mixing measure, whenever the displayed identity
and the almost-everywhere concentration condition make sense.  In
particular it excludes the conjectured, more restrictive decomposition
into extreme probability laws supported on transitive graphs.
\end{proof}

\paragraph{Formalization and verification.}
The accompanying \texttt{lean/Main.lean} uses actual Mathlib
\texttt{SimpleGraph}s, walk-defined induced balls, root-preserving graph
isomorphisms and their quotient, the cylinder-generated measurable
space, canonical measures of probability mass functions, and the
Lebesgue integral of nonnegative extended real functions.  The finite
graph component is the induced subgraph on actual reachable vertices;
its explicit rooted isomorphism to $P_3$ is proved.  The local law is
the pushforward of the uniform PMF on all $3(m+1)$ actual vertices,
and equality to $\mu$ is proved for every $m$.

The principal checked results are
\texttt{copies\_local\_weak\_limit},
\texttt{vertex\_count\_unbounded},
\texttt{obstruction\_excludes\_transitive}, and
\texttt{no\_transitive\_barycentric\_decomposition}.
No assertion that $\mu$ is extreme is required for the disproof, and
the proof does not challenge the usual ergodic decomposition theorem
without the false transitivity restriction.

The independent \texttt{verify.py} enumerates every simple graph on
at most six vertices, checks the radius-two degree obstruction at every
root, checks the $P_4$ boundary case, and computes the uniform rooted
laws of the first forty disjoint-union examples with exact rational
arithmetic.  These finite computations are consistency checks;
the arbitrary-graph and arbitrary-mixture statements are proved by Lean.

Lean 4.19.0 and Mathlib revision
\texttt{c44e0c8ee63ca166450922a373c7409c5d26b00b} are pinned.
In \texttt{lean/}, run:
\begin{verbatim}
lake exe cache get
lake build
lake env lean -DwarningAsError=true Main.lean
\end{verbatim}
Run \texttt{python verify.py} at the submission root.
Compile this paper with \texttt{tectonic main.tex}.
Actual fresh-build logs and hashes are included. Printed theorem
dependencies contain only standard foundational axioms.
The work received author, parent and internal cross-review;
no external independent review is claimed.
\begin{thebibliography}{9}
\bibitem{source} The Last Math Competition, Conjecture 00000009305.
Exact bilingual source included.
\bibitem{AL} D. Aldous and R. Lyons,
\emph{Processes on Unimodular Random Networks},
Electronic Journal of Probability 12 (2007), 1454--1508,
\href{https://arxiv.org/abs/math/0603062}{arXiv:math/0603062}, \S2.
\end{thebibliography}
\end{document}

